% Defect classes, particle classes and the charged-lepton register.
% Input by deriving_the_particle_zoo_from_observer_consistency.tex as its own section.
\section{Defect classes and particle representations}
\label{sec:defect-particle-classes}

The finite runs of the executable observer carry a protected register on the
faces of every carrier. A face holds a flux class $k\in\mathbb Z_6=\mathbb Z_3\times
\mathbb Z_2$, the product of a colour-centre charge and a weak-centre parity. Accepted
register twists create conjugate pairs of classes, the pairs move, fuse and
annihilate, and the classes on one closed carrier always add to zero modulo six. This
section compares these classes with Standard Model representation labels, states
the statistics under the declared quantum-code premises, and reads the native run
census. A face-flux class supplies no particle decoder, generation, chirality or mass.

\subsection{The class rule}

The internal space of one quantum is the exterior algebra of the five-level code
$V=\mathbb C^3_{-1/3}\oplus\mathbb C^2_{1/2}$ of Section~\ref{sec:finite-supply}. A
pattern has $n_3\in\{0,\dots,3\}$ colour levels and $n_2\in\{0,1,2\}$ weak levels
occupied, so there are $32$ patterns, and
\[
 6Y=3n_2-2n_3,\qquad Q_{\rm em}=T_3+Y .
\]

\begin{theorem}[Centre classes of the declared code]
\label{thm:defect-particle-classes}
On every one of the $32$ patterns the centre element
$(\omega I_3,-I_2,e^{2\pi iY})$ of $\mathrm{SU}(3)\times\mathrm{SU}(2)\times\mathrm U(1)$,
$\omega=e^{2\pi i/3}$, acts as the identity, and a flux of class $f$ acts by the
character $e^{2\pi i f\cdot 6Y/6}$. The class seen by a pattern is
\[
 k\equiv 6Y \pmod 6,\qquad k\bmod 3=n_3\bmod 3\ (\text{triality}),\qquad
 k\bmod 2=n_2\bmod 2\ (\text{weak centre}).
\]
The classes have multiplicities $(4,6,6,4,6,6)$ for $k=0,\dots,5$ and carry the
centre charges of the Standard Model representations in Table~\ref{tab:defect-classes}.
\end{theorem}

\begin{proof}
On a pattern with $n_3$ colour and $n_2$ weak levels, $\omega I_3$ acts by $\omega^{n_3}$,
$-I_2$ by $(-1)^{n_2}$ and $e^{2\pi iY}$ by $e^{i\pi(3n_2-2n_3)/3}=(-1)^{n_2}\omega^{-n_3}$; the
product is one. A flux of class $f$ acts by $e^{2\pi if\cdot 6Y/6}$, which depends on $6Y$ only
modulo six. Reducing $6Y=3n_2-2n_3$ modulo $3$ and $2$ gives the triality and the weak centre,
and counting patterns gives the multiplicities. The executable certificate checks all
$32$ cases in exact arithmetic \cite{ophrepo}.
\end{proof}

\begin{table}[htbp]
\centering\small
\caption{Defect classes and field representations with those centre charges.
The class is $6Y\bmod 6$; the table distinguishes the chiral representation labels
and includes their conjugates. A run class alone supplies no chirality or species assignment.}
\label{tab:defect-classes}
\begin{tabular}{@{}cllll@{}}
\toprule
$k$ & $\mathbb Z_3$ & $\mathbb Z_2$ & Field representations & Native class records\\
\midrule
$1$ & $\omega$ & $-1$ & left-handed quarks $Q=(u_L,d_L)$ & mobile; annihilates with $5$\\
$5$ & $\omega^2$ & $-1$ & their antiquarks $\bar Q$ & as class $1$\\
$4$ & $\omega$ & $+1$ & right-handed quarks $u_R,d_R$ & frozen after sweep $15$\\
$2$ & $\omega^2$ & $+1$ & their antiquarks $u^c,d^c$ & as class $4$\\
$3$ & $1$ & $-1$ & lepton doublet $L$, its conjugate, the Higgs doublet & frozen\\
$0$ & $1$ & $+1$ & $e_R,\nu_R$; photon, gluons, $W$, $Z$ & no defect\\
\bottomrule
\end{tabular}
\end{table}

The identification is exact at the level of centre classes. It decides which flux each
particle sees: a quantum of class $k$ taken around a flux of class $f$ acquires the
Aharonov--Bohm phase $e^{2\pi ifk/6}$. Colour flux ($f=2,4$) is seen only by quarks, weak
flux ($f=3$) only by the weak doublets, the left-handed fermions and the Higgs, and the
generating fluxes $f=1,5$ by every class except zero. The electrically charged $e_R$ sits
in class zero and sees none of these fluxes. The class carries neither the generation nor $T_3$, so $u$ and $d$, or
$e$ and $\nu$, share a class, as do a colour singlet and an octet. These finer labels sit
in the quantum's full pattern and in the generation register of
Section~\ref{sec:lepton-register}.

\subsection{Spin and statistics}

The binary icosahedral cover of the carrier rotations has one central involution $z$.
On the code it acts as
\[
 \rho(z)=\operatorname{diag}(1,1,1,-1,-1)=e^{i\pi\cdot 6Y},
\]
which equals $(-1)^{n_2}$ on every pattern and the holonomy of a class-three flux. On a
port space of flux $k$ the same element acts as $(-1)^k$; the weak levels sit at odd flux,
so the half-integer behind odd $6Y$ is the charge--monopole spin of the doublet levels.
The construction supplies a spin doublet on each quantum, the stated port-flux sectors
and antisymmetric occupations. On its declared Gauss-invariant sector the code,
register and link-end factors of $z$ multiply to one, and the rotation by $2\pi$ equals
$(-1)^{N}$, with $N$ the number of quanta.

\begin{proposition}[Statistics of the objects]
\label{prop:defect-statistics}
With the supplied spin-doublet, Gauss-sector and antisymmetric-occupation premises,
exactly one grading compatible with the declared vertices agrees with the $2\pi$
rotation. Under it every pattern is a fermion and the defects, the gauge links and
the Higgs register are bosons. Three quanta on one carrier have odd fermion parity.
\end{proposition}

The vertex constraints of the source alone admit sixteen gradings; on the declared Gauss sector four
survive, generated by the quantum-number parity and the coloured-quanta parity, and the
rotation selects one. The fermions are therefore the quanta whose pattern carries the class,
and the defect of that class is the flux the pattern sees. A carrier holding the pure flux
content $\{4,4,4\}$ has the centre classes of three quarks and the statistics of a boson; a
baryon is three fermionic quanta on the carrier. Antisymmetry of the occupations is forced
under the filled-reference hypothesis; this package does not construct that reference,
and the tick alone admits both occupation lifts.

\subsection{Reactions and closed carriers}

Every one of the $16$ reaction types of the runs conserves the class modulo six, and the
frequent ones carry the centre classes of the Standard Model vertices, with one diquark
channel
(Table~\ref{tab:defect-reactions}). Because a carrier is a closed screen, its classes add to
zero: a colour class never sits alone on a carrier.

\begin{table}[htbp]
\centering\small
\caption{Reaction channels of the level-10 run and the Standard Model vertex with the same
centre classes. Event counts are whole-run totals; the conjugate channel is in brackets.}
\label{tab:defect-reactions}
\begin{tabular}{@{}lrl@{}}
\toprule
Channel & Events & Vertex\\
\midrule
$1+5\to0$ & $609{,}494$ & quark--antiquark annihilation\\
$1+3\to4$ \ ($3+5\to2$) & $238{,}523$ \ ($239{,}773$) & quark mass vertex $\bar Q H q_R$\\
$1+2\to3$ \ ($4+5\to3$) & $315{,}024$ \ ($314{,}514$) & quark Yukawa vertex $Q\,d^cH^\dagger$\\
$0\to3+3$ & $89{,}909$ & lepton mass vertex $L\,H^\dagger e^c$\\
$1+1\to2$ \ ($5+5\to4$) & $458{,}209$ \ ($456{,}755$) & two quarks into an antitriplet\\
\bottomrule
\end{tabular}
\end{table}

The carrier contents of the level-10 run are exactly class-neutral: $\{2,4\}$ on
$69.16$ percent of occupied carriers, $\{1,5\}$ on $14.43$ percent, $\{3,3\}$ on $10.88$
percent, $\{2,2,4,4\}$ on $1.93$ percent, the triples $\{4,4,4\}$ and $\{2,2,2\}$ on
$1.20$ percent each, and rarer neutral mixtures such as $\{1,2,3\}$, the centre content of
the quark mass term. These contents share centre labels with possible meson,
baryon and lepton-pair representations. They do not identify those physical
objects; no carrier holds a net colour or weak centre charge.

\subsection{Population history of the finite runs}

The runs start from independent random loads and a blank register, so the preparation is
statistically homogeneous. In the level-10 run every class is created in pairs within the
first consistent slice: $124{,}786$ pairs $(1,5)$, $208{,}860$ pairs $(2,4)$ and $41{,}328$
pairs $(3,3)$. Classes $1$ and $5$ then annihilate ($609{,}494$ of $757{,}258$
pair births), while classes $2$, $4$ and $3$ freeze once the loads balance.
The conjugate-class counts are $3{,}053{,}295$ and $3{,}050{,}633$ for classes
$1$ and $5$. The conjugate triples have counts $86{,}476$ and $86{,}788$,
whose difference has magnitude $0.75$ times the square root of their sum.
This normalization is an arithmetic diagnostic; it supplies no independent
Poisson sampling law for correlated run histories or physical matter asymmetry.

Of all defect identities of the level-10 run, the classes $1,5,4,2,3$ take $24.1$, $24.1$,
$19.6$, $19.6$ and $12.7$ percent, and the shares at levels six and seven agree with these
within $0.3$ percentage points. A separate counting control assigns equal
weights to the fermionic code components and conditions on nonzero class.
Its 28 retained components give $6/28$, or $21.4$ percent, in each of classes
$1$, $2$, $4$ and $5$, and $4/28$, or $14.3$ percent, in class three.
The control is identical for the full and degree-restricted fermion menus
because their difference lies in class zero. It excludes the Higgs and
gauge components and supplies no thermal occupation law.

The run's pair-birth fractions are $50.6$, $43.4$ and $6.0$ percent for
$(1,5)$, $(2,4)$ and $(3,3)$. Folding the separate counting control into
these pair labels gives $42.9$, $42.9$ and $14.3$ percent. Component
multiplicities do not derive pair-production probabilities. The simulation
measures class populations and persistence on its declared event clock.
It supplies no physical temperature, species decoder, baryon number or
photon count, so it determines no baryon-to-photon asymmetry or early-Universe
composition. Section~\ref{sec:species-identifiability} states the exact
code-count inverse problem and its separate thermal comparison.

\subsection{Gauge group, couplings and lepton number}

The record law is covariant under independent phases on the twelve ports. Read on the
protected code, twelve relative port phases per seam generate exactly
\[
 G=\mathrm S(\mathrm U(3)\times\mathrm U(2))
 =\frac{\mathrm{SU}(3)\times\mathrm{SU}(2)\times\mathrm U(1)}{\mathbb Z_6},
\]
where $\mathbb Z_6$ is the kernel of the covering product. The centre of $G$ is the circle
$\{\operatorname{diag}(e^{2it}I_3,e^{-3it}I_2)\}$, and the defect register carries its
order-six subgroup generated by $e^{2\pi iY}$, which acts on a pattern by $e^{2\pi ik/6}$
with $k=6Y\bmod6$. The seam cost of a code state
carried around a plaquette, averaged over the code, is $1-\operatorname{Re}\operatorname{tr}_5
U_P/5$; this action gives equal bare couplings at its declared cutoff, with
$\sin^2\theta_W=3/8$ there, and the gauge waves travel on one cone with two polarizations
each \cite{ophGauge}. The native transports alone form a finite group; the continuous port
phases are the extension that makes the links continuous. Equality of bare couplings does
not establish equality of renormalized couplings at the scale of the comparison chain.

Every operation in the declared sector conserves
\[
 B-L=N-\tfrac23\,(N_{c1}+N_{c2}+N_{c3}),
\]
a combination of the number of quanta and the three colour-level numbers. It is anomaly free
on the sixteen states of one family with $\nu^c$, and it forbids number-violating
Majorana pairing in that sector. A Dirac neutrino mass additionally requires
the corresponding coupling and a physical species identification; conservation
of this charge alone supplies neither a neutrino mass nor that identification.

\subsection{The charged-lepton register}
\label{sec:lepton-register}

The proposed charged-lepton reading assigns three characters of a declared generation
register to the electron, muon and tau. The register is a construction on the coherent
seam law; a class-three defect alone carries no generation label. Its squared-eigenvalue
ratios can be compared with lepton masses, conditional on the stated preparation,
analyzer, clock and mass interpretation.

\paragraph{The repair event records its source.}
On the oriented address $(a,b)$ of a seam, with $p=(\ket a+\ket b)/\sqrt2$ and
$m=(\ket a-\ket b)/\sqrt2$, the coherent repair event is the swap of the port qubit $(p,m)$
with the record qubit. Read in its conjugate basis, the record holds the source endpoint
relative to the address. The oriented record $V$ changes the sign of record one at the
reversed address; $VQ$ equals the law with the screen address's unitary at both addresses, and
it keeps the repair branch that the plain record cancels.

\paragraph{The register.}
Two seam objects share one face $f$ of the carrier: object $C$ on a seam whose orientation is
the screen orientation and object $E$ on a second edge of the same hexagon around the $f$ axis.
Each has its load, record and address. A neutral register $F$ on the ports of $f$ carries the
third-turn $R=\rho_f$ of the face, with $R^3=I$. One tick is
\[
 U(\theta)=K_\theta\,T\,V\,(W_C\otimes W_E),\qquad K_\theta=e^{i\theta J},\qquad
 J=\lvert1\rangle\langle1\rvert_{r_C}\otimes\sigma_z(o_C),
\]
where $T$ applies $R^{\pm1}$ when $C$ repairs in orientation $\pm$ and $J$ is the oriented
current. Each object is prepared with activity $1/3$ and its load at the screen tail of its
seam. The analyzer $\Phi$ is the Fourier zero over the three flags $E$ agrees with $C$ idle,
$C$ repairs $+$ with $E$ idle, and $C$ repairs $-$ with $E$ idle.

\begin{theorem}[Exact compression to the balanced circulant]
\label{thm:lepton-register}
The full $432$-state tick is unitary, and its compression is
\[
 M(\theta)=\langle\Phi|U(\theta)|p\rangle=\frac{\sqrt2}{3}\,C(\theta),\qquad
 C(\theta)=\tfrac1{\sqrt6}\,I+\tfrac1{\sqrt{12}}\bigl(e^{i\theta}R+e^{-i\theta}R^{-1}\bigr).
\]
$M$ commutes with $R$, its eigenvalues on the characters $Rv_k=\omega^kv_k$ are
$\lambda_k=\frac{\sqrt2}{3}\bigl(\frac1{\sqrt6}+\frac1{\sqrt3}\cos(\theta+2\pi k/3)\bigr)$,
and $|b|/a=1/\sqrt2$. With masses proportional to $\lambda_k^2$,
\[
 Q=\frac{\sum_k m_k}{\bigl(\sum_k\sqrt{m_k}\bigr)^2}=\frac23
\]
for every $\theta$ in the positive chamber $|\theta|\le\pi/12$ modulo $2\pi/3$.
\end{theorem}

The balance is a count of branch weights. The $C$ object repairs in each orientation with
probability $1/12$, and the $E$ object holds with probability $1/6$, so the circulant
coefficients satisfy $a^2:b^2:b^2=1/6:1/12:1/12$. The phase is a count as well: the register's
exclusive read $X=[R_+]\oplus[H]$ has
\[
 P(X)=\tfrac1{12}+\tfrac16-2\cdot\tfrac1{72}=\tfrac29 .
\]

\paragraph{The clock.}
On every blank-record input with uniform address the seam program satisfies $\langle I-\operatorname{Re}U\rangle=
P(\text{record }1)$, so the energy clock reads read probabilities. Inside a tick, $N$ weak rounds
of the coupling $J\otimes X$ give the zero-return operator $A_N(t)$ with
\[
 \|A_N(t)-e^{-itP(X)J}\|\le t^2/N .
\]
The declared clock fixes the unit of $t$: since $I-S$ has eigenvalues $0$ and $2$, the edge
term $(\pi/4)(I-S)$ run for clock time $\tau$ reproduces the tick's root-swap gate exactly
when $\tau\in1+4\mathbb Z$, and the principal value is one unit per tick. At that unit the
oriented current carries the phase $\theta=P(X)=2/9$ per tick; at finite
$N$ the phase is $2/9-35/(2187N^2)+O(N^{-4})$. A dressing that commutes with $J$ leaves the
phase unchanged at every order. On this clock the family mode $k$ has the rate $\lambda_k^2$,
which is the mass reading of Theorem~\ref{thm:lepton-register}.

\begin{table}[htbp]
\centering\small
\caption{Squared-eigenvalue ratios of the declared register at $\theta=2/9$, compared
with charged-lepton masses. No measured mass enters the evaluation of the ratios;
the architecture and comparison are not blind. Reference values: CODATA for
$m_\mu/m_e$, $m_\tau=1776.93\pm0.09$ MeV.}
\label{tab:lepton-register}
\begin{tabular}{@{}lrrr@{}}
\toprule
& Register & Measured & Difference\\
\midrule
$m_\mu/m_e$ & $206.770316$ & $206.7682827(46)$ & $+9.83$ ppm\\
$m_\tau/m_e$ & $3477.4728$ & $3477.37\pm0.18$ & $+0.61\,\sigma$\\
$Q$ & $2/3$ & $0.666665$ & \\
\bottomrule
\end{tabular}
\end{table}

The proposed electron, muon and tau labels are assigned to the characters
$\omega^1$, $\omega^2$ and $\omega^0$ of the face third-turn; the half-turns that carry
$f$ to its opposite face reverse
the phase and swap the electron and muon labels, leaving the spectrum unchanged. The muon
residual of $+9.83$ ppm has a linear phase-displacement diagnostic of
$-1.75\times10^{-7}$. That displacement is not applied. A dressing commuting
with the oriented current leaves its phase unchanged. A physical correction
requires a derived mass-reading map and its coefficient; neither the
complete finite-channel error nor a selected radiative coefficient explains
the residual. The fixed comparison and its preparation budget are given
in Section~\ref{sec:complete-lepton-cycles}.

\paragraph{The register on the carrier.}
On the sixty-address carrier the count $2/9$ is the exclusive read of two cells with
independent address registers, read through the law's own records: in the defect's cell the
two signs of the native current are written with weight $1/12$ each, and one second cell
writes its hold with weight $1/6$. One register on one carrier gives only multiples of
$1/240$. A class-three flux keeps one two-load content, the pair on one port, and the kick of
the weak round is the commutator of an oriented toggle with the tick of that pair.
For two double-loaded cells the analyzer compression is $\sqrt{5/18}\,C(\theta)$,
whereas the specified native-pair analyzer gives $(\sqrt2/3)C(\theta)$.
The two readings have the same eigenvalue ratios and different absolute rates;
their traces $\operatorname{tr}M^\dagger M$ are $5/18$ and $2/9$, respectively.

\paragraph{The weak round over many ticks.}
One round is the tick of both cells, the kick $e^{-i(1/N)J\otimes X}$, the inverse tick, one
return test per cell and a reset after a negative test; $N$ rounds are one unit of clock time.
Every round reads the first-tick weights, and the phase per unit clock time is
\[
 \theta_N=N\arctan\frac{P\sin(1/N)}{1-P+P\cos(1/N)},\qquad P=\tfrac29,\qquad
 \theta_N=\tfrac29-\tfrac{35}{2187N^2}+O(N^{-4}).
\]
The register was executed on the sixty-address carrier for $10^4$ to $2\times10^6$ rounds. The
executed phases follow the closed form at every $N$ and at every elapsed number of rounds, and
their limit is $2/9$ within $10^{-13}$. The marked component of the state has squared
projections $4/81$, $11/324$, $125/1296$ and $55/1296$ onto the four return-test sectors;
these sum to $P$, not to one. For a current eigenvalue $j$, put $a=e^{-ij/N}-1$.
The actual four outcome probabilities are $|1+aP|^2$ and
$|a|^2(11/324,125/1296,55/1296)$, which sum to one.
With every outcome kept and the reset applied, the
coherence between two values $j,j'$ of the current is multiplied per round by
$1-P+Pe^{-i(j-j')/N}$, so the phase $2/9$ holds on the complete record and no branch is
selected. The forward tick alone does not keep the count: it reads $2/9$, $0.486$, $0.536$ and
$0.467$ at ticks one to four.

\paragraph{Rates and masses.}
The three characters of the face third-turn are the eigenvectors of the declared
native-pair reading. Their rates per unit clock time are
\[
 y_k=\lambda_k^2=\tfrac1{27}\bigl(1+\sqrt2\cos(\tfrac29+\tfrac{2\pi k}3)\bigr)^2,\qquad
 \textstyle\sum_ky_k=\tfrac29,
\]
and the executed rates agree with this form within $10^{-13}$. Table~\ref{tab:lepton-masses}
lists them. With the measured electron mass as a scale and the squared-rate interpretation,
the register gives the displayed conditional muon and tau values; the physical species and
mass attachment is an additional requirement.

\begin{table}[htbp]
\centering\small
\caption{Rest rates of the three characters at $\theta=2/9$ and the masses they give with the
electron mass as the only measured input. Measured values: CODATA and the Particle Data Group.}
\label{tab:lepton-masses}
\begin{tabular}{@{}llrrr@{}}
\toprule
Character & Lepton & Rate $y_k$ & Mass (MeV) & Measured (MeV)\\
\midrule
$\omega^0$ & $\tau$ & $0.2096935560$ & $1776.985$ & $1776.93\pm0.09$\\
$\omega^2$ & $\mu$ & $0.0124683656$ & $105.659414$ & $105.6583755(23)$\\
$\omega^1$ & $e$ & $0.0000603006$ & input & $0.51099895069$\\
\bottomrule
\end{tabular}
\end{table}

The register, its preparation and its analyzer are a construction on the coherent law, and the
executable certificate checks the compression, the counts, the round, the clock unit and the
carrier statements in exact arithmetic \cite{ophrepo}. Open inside it are four
items: the operation that brings two loads onto one port, the source event that applies the
kick, the rule that one partner is read per unit clock time, and the number $N_p$ of phase
rounds, which shifts $m_\mu/m_e$ by $-824/N_p$ at leading order. The balanced circulant was
first written down from the measured masses \cite{koide1983}; here its balance and its phase
are the register's own read counts. A mass in MeV is $m_k=F\,E_{\rm cell}\,y_k$ with one
number $F$; on the electroweak dictionary $F$ is fixed by the lepton Yukawa weight
$y_J\approx0.0108$, which is open.

\subsection{The integer lift of the register}
\label{sec:integer-lift}

Keeping the integer current of every applied twist unreduced lifts the $\mathbb Z_6$ register
to integers. On the runs at levels four to seven the lifted value reduces modulo six to the
stored class on every face of every stored cut, and the twenty face values of every carrier
sum to zero in all $68{,}515{,}840$ carrier states. As a set of values the lift takes $6Y$ of a
Standard Model multiplet on $99.7$ percent of the defect faces. It is no conserved label of
the member inside a class: the law reads the class only, a value of size four or more does
not travel with its defect, and one carrier content occurs with several integer words
\cite{ophrepo}.

\subsection{The coupling chain and the electroweak scale}
\label{sec:coupling-chain}

The gauge couplings and the electroweak scale follow one chain of five forms: the cell
energy $E_{\rm cell}=E_*P^{-1/2}$, the unification scale $M_U=E_*e^{-2\pi}P^{1/6}$, one-loop
running with slopes $(33/5,1,-3)$, the transmutation law
\[
 v=E_{\rm cell}\exp\!\Bigl(-\frac{2\pi}{4\alpha_U}\Bigr),
\]
and the pixel closure $\bar\ell_{\mathrm{SU}(2)}+\bar\ell_{\mathrm{SU}(3)}=P/4$ with heat times
$t_i=4\pi^2\alpha_i$ read at the electroweak anchor $\mu=m_Z(\mu)$. At the fixed point
$P=1.6309721$ the chain gives $1/\alpha_U=24.3166$ at $M_U=2.4736\times10^{16}$ GeV and the
values of Table~\ref{tab:coupling-chain}. These are tree values of the chain; the multiloop
pole comparison is Table~\ref{tab:twelve}.

\begin{table}[htbp]
\centering\small
\caption{Outputs of the coupling chain against measured values. The chain's forms are
declared; the comparison is made after the values were known.}
\label{tab:coupling-chain}
\begin{tabular}{@{}lrr@{}}
\toprule
Quantity & Chain & Measured\\
\midrule
$v$ (GeV) & $246.60$ & $246.22$\\
$m_Z$, $m_W$ tree (GeV) & $91.588$, $80.330$ & $91.188$, $80.369$\\
$1/\alpha_2(m_Z)$ & $29.605$ & $29.58$\\
$\alpha_s(m_Z)$ & $0.11834$ & $0.1180(9)$\\
$\sin^2\theta_W(m_Z)$ & $0.23074$ & $0.23122$\\
$1/\alpha_{\rm em}(m_Z)$ & $128.308$ & $127.95$\\
$1/\alpha(0)$ & $136.995$ & $137.036$\\
\bottomrule
\end{tabular}
\end{table}

Four parts of the chain have a source reading. The cell energy is one radian of phase per
tick on the cubic arrangement of carriers. The $4$ in the transmutation law is the inverse
plaquette cost of one class-three flux quantum: the central element
$\operatorname{diag}(1,1,1,-1,-1)$ has shortest generator $2\pi iT_3$ with
$\operatorname{tr}_5H^\dagger H=2\pi^2$, and the five nontrivial classes give the counts
$(12/5,3,4,3,12/5)$. The anchor is the first freeze of the weak currents: the class-three
register gives the neutral direction, which contains hypercharge, the largest mass, so the
collar law $d_R e^{-tC_2}$ with one coefficient for the three weak directions holds down to
$m_Z$, under one stated premise on the collar generator below the anchor. Hypercharge carries
no representation entropy because every irreducible representation of $\mathrm U(1)$ has
dimension one, and both factors $2\pi$ of the chain are the modular period.

The other parts are declared. The source as built has $16+\overline{16}$ patterns per copy,
twelve gauge currents and one complex Higgs doublet, with one-loop slopes $(81/10,5/6,-3)$;
on the same chain it gives $v=222.9$ GeV, $\alpha_s=0.1123$ and $\sin^2\theta_W=0.2056$. The
ratio $(1/\alpha_2-1/\alpha_3)/(1/\alpha_1-1/\alpha_2)$ at $m_Z$ equals $115/218$ for this
content and $5/7$ for the declared slopes; the measured couplings give $0.717$. The weak and
colour slopes of the declared chain equal those of the source's content with two doublets,
and the hypercharge slope does not. The stack's own unification scale is
$E_{\rm cell}e^{-2\pi}=1.7853\times10^{16}$ GeV, which gives $v=264.6$ GeV; the factor
$P^{2/3}$ between it and the declared form has no source reading. The displayed chain
uses collar weights $d_R$ and unit shape number as declared inputs. Its quoted numerical
precision requires control of higher-loop and threshold terms; one-loop evaluation alone
does not supply that control.

The transmutation law has the structure $v=E_{\rm cell}\,b\,\langle A\rangle$, where
$\langle A\rangle=\sum_jp_j(-1)^{2j}$ is the class balance of the seam link that the mass step
flips. In a leading weak-coupling form $\langle A\rangle\sim Me^{-cS}$, with
$S=2\pi/(4\alpha)$, the proposed law requires $c=1$ and $bM=1$. One free link in its one-tick thermal law has leading
weak-coupling exponent $c=1$, with prefactor tending to $2$ on the Cartan circle.
A link inside a lattice shares the half turn with its
surroundings: in the infinite-volume Gaussian limit on the Cartan circle,
$c=1-2/d$ on a cubic lattice in $d$ dimensions, and Monte
Carlo for $\mathrm{SU}(2)$ with the heat-kernel weight gives $c=0.5001(8)$ on $6^4$,
$0.3347(12)$ on $8^3$ and $0.3988(5)$ on the icosahedral space. None of the ensembles
measured here has coefficient one; the single-link value requires a pinned plaquette.
The sector law that pins
the seam link is open \cite{ophrepo}.

\subsection{Lattice ensembles and thermal comparisons}
\label{sec:early-universe-lattice}

The stored runs measure finite repair counts. The thermal calculations below use separate
Euclidean lattice ensembles, with temperature $T=1/(aN_t)$ in units
$\hbar=c=k_B=1$. The lattice spacing $a$ and its physical calibration are inputs;
Euclidean time slices are not identified with the repair clock. These ensemble
comparisons do not supply a shared cosmological cooling history.

\paragraph{The electroweak crossover.}
With the couplings of the declared chain as inputs and the published reduction to the
three-dimensional $\mathrm{SU}(2)$ and Higgs theory \cite{pwklrs1996reduction}, the transition
has crossover central value
\[
 T_c=161.1298\ \mathrm{GeV}.
\]
The Standard Model lattice value is $159.5\pm1.5$ GeV \cite{pwdonofrio2016crossover}.
The same pipeline on Standard Model inputs gives
$159.76$ GeV against the published $159.58(6)$ GeV. On the chain the first-order line ends at
a Higgs mass of $72.5\pm0.7$ GeV. With the assumed Standard Model thermal census
$g_*=106.75$ and a radiation-dominated Friedmann relation, the crossover corresponds to
$9.0\times10^{-12}$ s. Converting this to $1.3\times10^{32}$ repair ticks additionally
assumes $\hbar/E_{\rm cell}=6.885\times10^{-44}$ s per tick. With the declared source
content and assumed Higgs and top masses the
crossover lies between $145$ and $149$ GeV.

The reduction uses $y_c=y(T_{\rm lat})+y_{\rm fit}-y_{\rm ref}$;
the reference ensemble enters all three terms. Paired peak replicas are
absent from the retained data, so the statistical standard deviation is
not identified. Linear propagation of the supplied marginal variances,
with independence between distinct lattice sets and arbitrary covariance
within the reference set, gives the sharp compatible range
$[0.11810,0.13680]$ GeV. Without the cross-set independence premise the
range is $[0,0.27631]$ GeV. These are uncertainty ranges under a fixed
linearized fit, not confidence intervals. The declared one-percent
reduction allowance is $1.6113$ GeV and has no supplied probability law;
it is not combined in quadrature with the statistical range or used to
quote a standardized agreement. A separate comparison at matched physical size
$Lg_3^2=8$ gives $160.8546\pm0.0720$ GeV for a linear cutoff fit.
Linear, quadratic and omitted-set variants span a maximum displacement
of $0.3490$ GeV from that value. This finite-design sensitivity is separate
from the reduction estimate and does not bound omitted cutoff powers or
the infinite-volume limit. Three points fitted to a constant, linear and
quadratic term have zero residual degrees of freedom. The measured
susceptibility-volume exponent decreases from $0.227(38)$ to $0.089(89)$
across the two largest-size comparisons; the first-order control grows
from $2.163(32)$ to $2.881(88)$, toward the volume exponent three.

\paragraph{Colour deconfines.}
Restricted to a colour holonomy $W$, the seam cost $1-\operatorname{Re}\operatorname{tr}_5U/5$
is $(3/5)(1-\operatorname{Re}\operatorname{tr}W/3)$, the Wilson plaquette action with
$\beta=6/g^2$. The colour centre $\mathbb Z_3$ is the flux classes $2$ and $4$ of the register,
and the classes charged under it are the quark classes $1,2,4,5$. A Monte Carlo study of these
links at $N_t=4$ gives the transition coupling $\beta_c=5.69105(39)$ from three volumes and
$5.69248(71)$ from the two largest, against the published $5.6925(2)$
\cite{pwboyd1996thermodynamics}. The susceptibility maximum grows as $N_s^{2.555(29)}$, the
histogram of the Polyakov loop has two peaks, and the deconfined runs settle in the three
$\mathbb Z_3$ sectors: the transition is of first order. With the string tension from the same
lattice,
\[
 T_c/\sqrt\sigma=0.614(18),
\]
against the published $0.598(3)$ at this lattice spacing and $0.629(3)$ in the continuum
limit. The study holds the links alone at one lattice spacing and scans the coupling as an
ensemble parameter; continuous links on the seams are a declared step.

\paragraph{Quarks bind.}
Quenched Wilson quarks on the same links at $\beta=5.7$ on $12^3\times24$, with $128$
configurations, give a static potential that rises linearly, $a^2\sigma=0.1667(17)(135)$,
and in the chiral limit
\[
 m_N/m_\rho=1.45\pm0.13,\qquad m_\rho/\sqrt\sigma=1.333(122),\qquad m_N/\sqrt\sigma=1.932(182).
\]
The normalized-mass errors combine same-stream jackknife errors with the largest displacement
over crossed hadron and potential fit choices. They do not bound quenching, finite-volume
or continuum effects.
Keeping the paired jackknife replicas also determines the joint mass
covariance. The same-stream contrast is
$(m_N-m_\rho)/\sqrt\sigma=0.59903\pm0.04286$ statistically, whereas
discarding the covariance gives $0.04407$ for that statistical error.
This contrast does not include fit-systematic variation. The common scale
cancels from $m_N/m_\rho$ in every paired replica. The retained covariance
therefore belongs in joint comparisons of the normalized masses.
The measured ratio $m_N/m_\rho$ is $1.211$. Above the transition the binding melts: the slope
of the free energy of a static pair falls from $0.9896(51)$ per lattice unit at $\beta=5.6$
to $0.00107(1)$ at $\beta=5.85$. Masses in GeV need the light-quark masses and the lattice
spacing in GeV, which are open. With the chain's $\alpha_s(m_Z)$ and imported lattice ratios
the colour scale is $\Lambda^{(3)}=0.340$ GeV, the nucleon $0.933(47)$ GeV and the rho
$0.773(42)$ GeV.

\paragraph{Standing of the three comparisons.}
The colour and hadron numbers are those of pure and quenched $\mathrm{SU}(3)$ with the Wilson
action. Their agreement with published lattice values confirms that the model's colour links,
once continuous, are that theory. The crossover temperature takes the declared chain's
couplings as inputs. The class shares of the stored runs lie near equal occupation per state
without equalling it, and the baryon excess and the light-element abundances are not
computed.

\subsection{Complete cycles and controlled lepton rates}
\label{sec:complete-lepton-cycles}

The finite lepton instrument can be composed without discarding a return failure.
Prepare each cell's sixty-address register in the uniform state $u$, and let
$S=u\otimes u$. A Householder unitary gives an explicit map from a blank address
basis vector to $u$. The blank double loads, independent fresh cells and this
preparation service are declared inputs. After the native forward tick, the
exclusive-read kick and the inverse tick, the joint state has the form
\[
 S\otimes\psi+X\otimes(U_h-I)\psi,
 \qquad U_h=e^{-ihJ},\quad J=\operatorname{diag}(-1,0,1),
 \quad \langle S,X\rangle=\|X\|^2=p=\frac29.
\]
The two return tests have four outcomes. Their complete Kraus operators are
\[
 A_{11}=(1-p)I+pU_h,\quad
 A_{01}=\sqrt{\frac{11}{324}}(U_h-I),\quad
 A_{10}=\sqrt{\frac{125}{1296}}(U_h-I),\quad
 A_{00}=\sqrt{\frac{55}{1296}}(U_h-I).
\]
The first outcome includes interference with the standing source. The sum of
the four effects is identity, and the unconditional register channel is
\[
 \Phi_h(\rho)=(1-p)\rho+pU_h\rho U_h^*.
\]
Exchanging the old source pair with an independently prepared fresh pair resets
the active apparatus. The old pair, its flag and all its correlations are
exported. This unitary exchange supplies the reset while retaining every
outcome; it does not erase an unknown quantum state. It requires fresh storage
and a preparation independent of the continuing register and reference.

The mass-read analyzer gives a separate complete channel. In the three Fourier
characters its compressed effect is $E=\operatorname{diag}(e_k)$, where
\[
 e_k=\left[\sqrt{\frac5{108}}+
       2\sqrt{\frac5{216}}\cos\left(\theta+\frac{2\pi k}{3}\right)\right]^2,
 \qquad \operatorname{tr}E=\frac5{18}.
\]
The pulled-back analyzer vectors $X_k$ satisfy
$\langle X_l,X_k\rangle=\delta_{kl}e_k$, by the equal three-character weights
and $1+\omega+\omega^2=0$. The sum of their nonreturn Gram matrices is
$\operatorname{diag}(e)-ee^{\mathsf T}$, positive by weighted Cauchy--Schwarz.
It has rank three at $\theta=2/9$, so a single scalar failure operator does
not reproduce the source environment. With $z=e^{-ih}$, the complete mass
channel fixes populations and, for $k\ne l$, obeys
\[
 \Psi_h(\rho)_{kl}
   =\bigl[1+(z-1)e_k+(\bar z-1)e_l\bigr]\rho_{kl}.
\]
The double-load rates have common factor $5/4$ relative to the native-pair
rates of total $2/9$. Their ratios agree; their clock normalization is explicit.

For $N$ independent fresh-source cycles with $h=t/N$, finite channel estimates
give
\begin{align*}
 \frac12\|\Phi_{t/N}^{N}-\operatorname{Ad}(e^{-itpJ})\|_\diamond
   &\le \frac{14t^2}{81N},\\
 \frac12\|\Psi_{t/N}^{N}-\operatorname{Ad}(e^{-itE})\|_\diamond
   &\le \frac{349t^2}{1296N}.
\end{align*}
The first estimate centers the binomial kick count and uses the spectral
diameter $D(J)=2$. The second compares the Taylor remainders of the source
projector, of diameter one, and its compression, whose diameter is at most
$5/18$. Channel contraction and telescoping extend both estimates to arbitrary
untouched references. At $t=1$, the conservative ideal cycle counts for error
at most $10^{-6}$ are $172{,}840$ and $269{,}291$. Per-cycle channel error
$\epsilon$ and independent preparation error $\eta$ add at most
$N(\epsilon+\eta)$. Marginal preparation fidelity alone does not bound
correlations with the register. Diagonal densities are stationary, including
a faithful uniform density; this does not select a unique reference. The
exported history grows. No physical duration or mass-pole interpretation is
provided by these finite estimates.

The spectrum's parameter dependence can be calculated separately. For the
shape $1+c\cos(\theta+2\pi k/3)$, the positive chamber gives
$Q=1/3+c^2/6$. At $(\theta,c)=(2/9,\sqrt2)$, the Jacobian of the logarithmic
muon/electron and tau/electron ratios is
\[
 \begin{pmatrix}56.14388&32.61133\\51.22679&34.45440\end{pmatrix},
 \qquad \det=263.82944.
\]
The two shape parameters are locally identifiable from the two ratios; the
common normalization cancels. The CODATA 2022 comparison value
$206.7682827\pm0.0000046$ leaves a muon/electron residual of $+9.83358$
parts per million, or $442.02$ experimental standard uncertainties. No parameter
or round count is adjusted to cancel it. Fixed individual flag marginals
$1/12$ and $1/6$ permit an exclusive phase anywhere in $[1/12,1/4]$;
independence fixes their joint probability to $1/72$ and the phase to $2/9$.
Thus the full preparation, rather than the marginals alone, supplies that
phase premise. The executable source arrays, complete instruments and
finite-history checks are in \texttt{model/lepton\_cycle} \cite{ophrepo}.

The same fixed shape gives $m_\tau/m_e=3477.47283710$. Dividing the
Particle Data Group 2024 average $m_\tau c^2=1776.93\pm0.09$ MeV by the
CODATA 2022 electron energy $0.51099895069\pm0.00000000016$ MeV gives
$3477.36526191$ \cite{pwnavas2024review}. The electron and direct muon-ratio
inputs are listed in the \href{https://physics.nist.gov/cuu/Constants/Table/allascii.txt}{2022 CODATA table}.
The selected tables supply no cross-edition covariance. For $r=t/e$, put
$a=u_t/e$ and $b=t u_e/e^2$. All positive-semidefinite covariances give
$|a-b|\le u_r\le a+b$ under linear propagation. Here this range is
$[0.17612452,0.17612670]$, so the tau residual is $30.93584$ parts per
million or $[0.610783,0.610791]$ experimental standard uncertainties.
These reference-only discrepancies supply neither a joint significance nor
a physical-theory error. The mass-pole interpretation is unproved.
The comparison is retrospective; measured masses and Koide relations were
known during architecture development. The executable comparison freezes
both ratios and selects its finite round counts from a channel-accuracy
target independent of the mass residuals.

Preparation uncertainty has a separate, finite statement. For a fixed
analyzer projector and independent source states at trace distance at most
$\eta$, the compressed effects differ in operator norm by at most $\eta$.
This follows by applying the variational trace-distance bound to the source
effect induced by each unit register vector. Weyl's inequality bounds each
ordered rate by $\eta$. If $\eta<e_e$ and twice $\eta$ is below all spectral
separations, the mode labels are preserved and
\[
 \frac{e_j-\eta}{e_e+\eta}\le\frac{e'_j}{e'_e}
 \le\frac{e_j+\eta}{e_e-\eta}.
\]
These are fixed-analyzer effect-rate bounds. Complete-channel distance and
source preparation errors are not adjustable physical mass corrections.

\subsection{An exact finite population and its symmetry}
\label{sec:finite-population-dynamics}

A closed triangular cell supplies a finite test of persistent populations.
It has three loads in $\{0,\ldots,5\}$ and two oppositely oriented face
classes $(k,-k)$ in $\mathbb Z_6$. Each attempt selects one of the three
oriented seams uniformly and a fair balancing coin. It uses the same native
load repair and class-acceptance functions as the federation. The state space
has $1296$ elements and the complete instrument has $7776$ input-event
branches, including load-only and no-operation events. The geometry is a
declared two-face cell; no transfer of its population to the full federation
is asserted.

With class cost $c=(0,5,8,10,8,5)$ and $\kappa=5$, the quantity
\[
 V(x,k)=\kappa\sum_{a=0}^{2}x_a^2+2\bigl[c(k)+c(-k)\bigr]
\]
does not increase on any branch. An accepted current $j$ changes the load
square by $-2j(d-j)$ and the register cost by at most $\kappa j(d-j)$.
Rejected twists retain only the nonpositive load contribution. The complete
kernel conserves total load and total face class. Its finite harmonic
absorption equations determine the persistent occupied-face probability
from every initial state. For independent uniform initial loads and zero
initial flux, their exact average is
\[
 \Pr(\text{persistent occupied faces})=\frac{1007}{13824}
       =0.0728443287037\ldots.
\]
The full distribution after 128 attempts bounds this limiting share in an interval of width less than
$5.86\times10^{-10}$. These probabilities describe operational repair events,
without a physical temperature, cooling law or species-abundance dictionary.
For every initial distribution on all $1296$ states, the probability of
remaining outside the terminal set after $n$ attempts is at most
$25^{-\lfloor n/32\rfloor}$, from the exact all-state 32-attempt
absorption bound.

The native law also has an exact conjugation. Send every load $x$ to $5-x$,
every face class and integer lift to its negative, and every coin to its
complement. The repaired current changes from $j$ to $-j$ while the release
$j(d-j)$ and the class-cost comparison are unchanged. Exhaustive enumeration
checks all $2592$ local load/class/coin branches; the same identity verifies
all $7776$ branches of the population cell. It extends pathwise on the closed
federation with an equivariant seam schedule.

For a finite kernel $K$ and conjugation $C$, the conditions
$K(Cx,Cy)=K(x,y)$ and $\mu(Cx)=\mu(x)$ imply
$(\mu K^n)(Cx)=(\mu K^n)(x)$ for every finite $n$.
Every odd observable $f(Cx)=-f(x)$ then has zero expectation. These finite-sum
identities are formalized in the Lean dynamics library; the native branch
checker supplies the source-specific premises. The signed class differences
are odd, though single histories can fluctuate. If the initial symmetry
defect is $a_0$ in total variation and the complete kernel at attempt $r$
has equivariance defect at most $\varepsilon_r$, maximal coupling gives
\[
 |\mathbb E f|\le
 1-(1-a_0)\prod_{r=1}^{n}(1-\varepsilon_r),\qquad |f|\le1.
\]
A coin of probability $p_r$ contributes at most $|2p_r-1|$ to that defect.
Conjugation-even selection of probability $q$ can amplify the bound by at
most $1/q$; asymmetric selection can produce a conditional imbalance from
an exactly symmetric law. The complete records therefore distinguish
dynamical asymmetry from selection. Class conjugation alone does not identify
baryon number or predict the baryon-to-photon ratio. Reproduction is supplied
by \texttt{model/native\_abundance\_epochs} and
\texttt{model/charge\_conjugation\_dynamics} \cite{ophrepo}.

\subsection{Species information in a centre-class census}
\label{sec:species-identifiability}

For a declared five-bit exterior-code component, let $n_3$ and $n_2$
count colour and weak occupations. Its centre class is
$k=(3n_2-2n_3)\bmod6$. The matrix $A_{kj}=1_{k(j)=k}$ pushes component
populations $x\ge0$ to the six-class census $c=Ax$. The full 32-component
menu has multiplicities $(4,6,6,4,6,6)$, rank six and nullity 26.
Removing exterior degrees zero and five, the two right-neutrino labels,
gives 30 components and nullity 24. Adding four Higgs and twelve gauge
internal components gives nullities 42 and 40, respectively; physical
helicity multiplicities are a separate count.

An aggregate $w\cdot x$ is identifiable from every census precisely when
$w$ is constant within each class. Necessity follows by exchanging two unit
populations in that class; sufficiency follows by summing the shared
coefficient against $c_k$. More generally the sharp nonnegative bounds are
\[
 \sum_k c_k\min_{j:k(j)=k}w_j\ \le w\cdot x\le
 \sum_k c_k\max_{j:k(j)=k}w_j.
\]
Assigning all class population to a minimizing or maximizing component
attains either endpoint. Thus the quark-labelled total
$q=c_1+c_2+c_4+c_5$ is identified on this declared component dictionary.
With bosons admitted, the lepton total spans $[0,c_0+c_3]$, the gauge
total $[0,c_0]$ and the Higgs total $[0,c_3]$. These are separate marginal
bounds; their endpoints need not occur together.

One additional fermion/boson bit within classes zero and three suffices
for all four totals. To distinguish also up/down and charged/neutral
lepton aggregates requires three label values in the boson extension:
the electron, neutrino and Higgs signatures in class three prove the lower
bound. Class-dependent signature labels attain it. This is a minimal
information theorem for the declared dictionary, not a physical decoder.
Native face-flux events, including zero flux, do not by themselves count
code quanta or thermal particles.

Uniform component weights give quark/lepton counts $24:8$ in the full
fermion menu and $24:6$ in the restricted one. They supply no occupation
probability law. For an explicit comparison, a massless ideal Standard
Model gas at one common temperature of $1$ TeV, zero chemical potentials,
three active-neutrino families and no populated right-handed neutrinos has
90 fermionic and 28 bosonic helicity states. Number densities weight
fermions by $3/4$ relative to bosons; energy densities use $7/8$
\cite{pwnavas2024review}. This ideal-gas benchmark uses the thermodynamic
conventions of \href{https://doi.org/10.1088/1475-7516/2018/05/035}{Saikawa and Shirai (2018)};
it omits their interacting equation-of-state corrections.
The effective totals are therefore $191/2$ and
$427/4$. Quark and lepton fractions among all particles are $108/191$ and
$27/191$ by number, versus $252/427$ and $63/427$ by energy. Their
fermion-only split is $80\%:20\%$. These are ideal-gas reference values
with interactions, masses and thermal-mass corrections omitted. No run
temperature, equilibration law or early-Universe abundance match is
inferred. The rational maps, extremizers, minimum labels and thermal
counting receipt are in \texttt{model/native\_abundance\_epochs} \cite{ophrepo}.

\subsection{Scalar response of localized defect modes}
\label{sec:defect-scalar-overlap}

A Yukawa coefficient can be defined operationally through the response of
a fermionic creation--propagation--detection amplitude to a scalar
intervention. In canonical field coordinates, the convention
\[
 M(h)=M(0)+\frac{h}{\sqrt2}Y_{\mathrm{eff}}+O(h^2)
\]
defines the zero-scalar-momentum mass-kernel vertex. Its derivation requires
the source intervention and field normalizations. Charge and topology
constrain allowed tensor channels, whereas their amplitudes depend on the
mode profiles and dynamics. This definition supplies no independent
interaction coefficient when the complete source specifies its
response. An on-shell particle interpretation requires the associated
propagation and state.

There is a finite overlap calculation on the specified compact texture
with cores at $(0,0,\pm1)$, texture amplitude $M=64$, and radius-three
exterior box. Its complete 686-row Hamiltonian has four isolated modes.
Let $W$ be an orthonormal frame for their entire spectral subspace and
$C=D_{\mathrm{stagger}}i\sigma_2$ its particle--hole pairing map. The
source identities give
\[
 C\overline W=WS,\qquad
 S=W^\dagger C\overline W,\qquad S^T=-S,
 \qquad S^\dagger S=I_4.
\]
The unit local scalar insertion at site $x$ has pairing kernel $P_xC$.
Its complete band response is
\[
 B_x=W^\dagger P_xW,\qquad
 Y_x=W^\dagger P_xC\overline W=B_xS.
\]
The same spatial factor can multiply a unit symmetric internal Higgs
tensor. No family coefficient is inserted in this spatial calculation.
All 343 site kernels are retained. A change $W\mapsto WU$ gives
$Y_x\mapsto U^\dagger Y_x\overline U$, so the singular spectra are
independent of a basis choice inside the band.

For a complete orthonormal complement $Q$, define
$L_x=Q^\dagger P_xC\overline W$. Every cross-site leakage obeys
\[
 L_x^\dagger L_y
   =S^\dagger(\delta_{xy}B_x-B_xB_y)S.
\]
Thus a real scalar profile $f$ gives $Y[f]=\sum_x f_xY_x$ with
\[
 L[f]^\dagger L[f]
  =S^\dagger\left(\sum_x f_x^2B_x-B[f]^2\right)S,
 \qquad B[f]=\sum_x f_xB_x.
\]
The uniform profile has $Y[1]=S$, four unit singular values, and zero
leakage. A uniform interaction on this carrier supplies no spatial
hierarchy.

For the specified 27-site intervention
$f_x=1_{\|x-(0,0,1)\|_\infty\le1}$, the numerical singular values are
\[
 (0.999999981366,\ 0.999999981366,\
  0.000061031430055,\ 0.000061031430055).
\]
The largest-to-smallest ratio is approximately $16385$. The weaker
channel reaches the patch through the opposite core's spatial tail.
The first-hop transfer Gram is $I_2/4$, yielding leading strength
$1/(4M^2)$. Complete controls at $M=32,64,128$ give
$M^2y_{\mathrm{weak}}=0.2499389,0.2499847,0.2499962$.
These values are numerical diagnostics, with no outward error interval
asserted for the ratio. No measured mass selects the profile or source
parameters.

The exterior resolvent supplies the normalization explicitly. In core and
exterior coordinates, $h=\left(\begin{smallmatrix}0&B^\dagger\\B&D\end{smallmatrix}\right)$
and an actual band eigenvector is
\[
 \psi=(u,-(D-E)^{-1}Bu),\qquad
 \|\psi\|^2=u^\dagger\bigl[I+B^\dagger(D-E)^{-2}B\bigr]u.
\]
Every band mode is reconstructed with this norm. Complete 1372-row Nambu
interventions also retain their complementary Schur response; the
nonzero spatial band Hamiltonian is kept. Dividing $f$ by its Euclidean
norm rescales every strength together. This coordinate normalization does
not determine a physical Higgs kinetic residue.

The result exhibits two different scalar-response strengths, each repeated
twice, for modes carrying the same internal current representation.
A dynamical scalar Hamiltonian must determine which contraction $Y[f]$
describes its collective mode. The specified profile is an intervention,
and the four spatial modes have no three-generation assignment. Physical
Yukawa values and relativistic masses are not inferred from these finite
overlaps. The complete site kernels, complement identities and spectral
controls are reproduced by \texttt{model/defect\_mode\_overlap} \cite{ophrepo}.

\subsection{Dynamical scalar selection and native contact response}

On the same compact texture, declare the scalar stiffness and its positive
frequency operator
\[
 K=I+L_{\mathrm{graph}}+M^2\operatorname{diag}(|n_x|^2),
 \qquad \Omega=\sqrt K.
\]
The scalar carrier consists of its blank and all site-labelled particle
and antiparticle doublets, with one global excitation cap. Its dimension
is $1+4(343)=1373$. Each of the 343 spatial scalar modes is retained.
For the scalar blank, the literal finite matrix units give
\[
 D_\phi(z)=(z-\Omega)^{-1}-(z+\Omega)^{-1}
          =2\Omega(z^2-K)^{-1}.
\]
Consequently $q=(2\Omega)^{-1/2}\phi$ has response $(z^2-K)^{-1}$.
A unit coupling to a ladder coordinate $\phi_j$ becomes
$\sqrt{2\omega_j}$ times that coupling in the coordinate $q_j$.
This normalization belongs to the blank reference; an interacting ground
requires its own response.

The lowest scalar mode is unique, positive and even under exchange of the
two cores. Its stiffness eigenvalue is approximately $6.99829101533$.
Its two doubled overlap strengths are $0.706890884$ and $0.706804630$,
with ratio $1.000122033$. Thus the large hierarchy of the specified local
probe is absent for this selected scalar mode.

The interaction can also select a scalar response without imposing a
classical profile. Fill every negative spatial fermion mode with all
fifteen native labels and take the scalar blank. The resulting Slater sea
is a global gauge singlet because all twelve native generators are
traceless. Retain its conjugate sea as an idle reference. Let $P$ be the
complete negative spectral projector and define the spatial pair Gram
\[
 G_{xy}=D_xD_y\det(P_{xy}).
\]
The full unit up, down and lepton tensors have exact internal Gram
$\operatorname{diag}(6,6,2)$. Wick contraction gives
$\|H_Y|s\rangle\|^2=56\operatorname{Tr}G\simeq29.01797916$.
The complete scalar mode probabilities in this first interaction image
put approximately $0.964506587$ in the odd bound mode and $0.000411883$
in the even lowest mode. The dominant reduced response profile has four
equal spatial strengths near $0.706847830$. These are first-image records,
not a dressed Higgs pole or a selected vacuum expectation.

The orthonormal pair $|s\rangle$, $H_Y|s\rangle/\|H_Y|s\rangle\|$
gives a two-dimensional Rayleigh--Ritz matrix for the full law, retaining
all native channels and scalar states. The numerical upper-bound estimate
is $E_{\mathrm{ground}}-E_{\mathrm{sea}}\le-2.46785318$. Independently,
the complete certified spatial gap and operator bounds prove the strict
rational inequality
\[
 E_{\mathrm{ground}}\le E_{\mathrm{sea}}
 -\frac{1456224000}{169057134218809}.
\]
The ground is taken in the same global singlet domain. The full sea is
therefore not its ground. A complete spectral instrument has success
operator $H_Y/(28\cdot343)$ and failure effect
$I-H_Y^2/(28\cdot343)^2$; every scalar mode and the original failure
probability are retained. These statements are reproduced by
\texttt{model/defect\_scalar\_modes} \cite{ophrepo}.

For the one-site texture-free native control, all thirty fermion modes
and the five-state scalar are retained. The equation-of-motion insertion
\[
 \mathcal K_{ab}=\{[c_a,H],c_b\}
\]
is the exact contact pairing derivative. Contraction with the actual
scalar Gram reconstructs all its components from scalar--fermion
three-point records. The lepton-only 320-state control has scalar first
moment $Z_H=14/5$. With all quark vertices active, the local native-current singlet ground is
different. Writing $\Delta=\sqrt{1+4\lambda}$,
$e=(1-\Delta)/2$ and $n=(1-1/\Delta)/2$, it has
\[
 Z_H=2-\frac74 n-\frac34 e\simeq9.58523463744.
\]
Here $\lambda$ is the selected source root, approximately $137.35663135$.
All 900 fermion pairs give 56 nonzero oriented contact components.
The three unit native channels have canonical contact modulus
$\sqrt{Z_H}\simeq3.09600301$. Their complete scalar-pulse responses differ:
the normalized pair-pointer slopes are approximately $-0.2108914572$
for either quark channel and $-0.1983104425$ for the charged lepton.
This is a state-specific contact and intervention result, without an
on-shell Yukawa or observed family assignment. The corresponding exact
spectral records and finite pulse controls are given by
\texttt{model/defect\_family\_vertices} \cite{ophrepo}.

The same local source also has a complete Higgs-dressed neutral-composite
spectrum, despite its elementary neutral contact nullspace. All nine target
roots are retained on the 26-dimensional target in each of four
spin/addition/removal settings. The resulting 44 original records comprise
36 success energies, four failures and four zero-weight outside outcomes.
The minimum success energy lies in the exact interval
\[
 \left[\frac{6301896140017}{70368744177664},
       \frac{25207584560069}{281474976710656}\right],
\]
approximately $0.089555330476$ source units. This gapped local composite
has no physical neutrino species or dispersion assignment.

The common finite construction has full CAR action on thirty native
fermion modes per site and retains the global scalar cap. Its texture-free
three-site member executes a scalar preparation and pair detection on all
ninety fermion modes. The literal supports of $H^n\psi$, for $n=0,1,2,3$,
have $1,11,664,7799$ coordinates, with coincident amplitudes combined in
exact arithmetic. If $P_R$ detects a complete right-site pair, then
\[
 P_R\psi=P_RH\psi=0,\qquad
 \|P_RH^2\psi\|^2=8(2/3-\sqrt2/2)^2>0.
\]
The spectral Taylor remainder is bounded by
$t^3\|H^3\psi\|/6$. At $t=10^{-4}$ this gives a rigorous receiver
probability between approximately $2.70\times10^{-19}$ and
$3.90\times10^{-19}$, with exact radical endpoints and the complete
complementary outcome retained. The dual charged preparation reference
is retained throughout. This is a finite interaction event, without a
particle mass or finite-speed interpretation; $\sqrt K$ directly couples
the outer scalar sites. The full source columns and error bounds are
reproduced by \texttt{model/defect\_mass\_join} \cite{ophrepo}.

\subsection{Propagation tests and a common energy unit}

A periodic extension of the compact texture gives exact supercell momentum
labels. For even cell side $L\ge6$, the hopping norm is at most three and
the on-site spectrum has four zero entries and all other entries $\pm M$.
Ordered-eigenvalue bounds therefore give exactly four bands in $[-3,3]$
and a complement with $|E|\ge M-3$, at every Bloch phase when $M\ge8$.
The complete Bloch spectra for $L=6,8$ and $M=8,64$ are retained. Direct
construction of a two-cell ring verifies the Fourier intertwiner on all
rows. Filled-sea additions and removals, each scaled by the inverse square
root of the one-body dimension, give a complete energy and crystal-momentum
instrument by the canonical anticommutation identity.

The low bands are nearly flat and depend strongly on direction. In the
$L=6$, $M=8$ control, positive-band secants of $E^2$ against $p_z^2$ are
negative between Bloch phases zero and $\pi/2$ along the third axis.
These numerical controls do not
establish relativistic rest masses from the zero-momentum gaps. The
periodic geometry also differs from an isolated moving defect and from
the interacting source. The construction and complete complementary
spectra are reproduced by \texttt{model/defect\_dispersion} \cite{ophrepo}.

For a fixed source clock, a GeV conversion multiplies every energy by one
positive constant $a$. A positive source interval $[x_-,x_+]$ and positive target
interval $[y_-,y_+]$ permit exactly
\[
 a\in[y_-/x_+,\ y_+/x_-].
\]
A common scale exists for independent interval anchors precisely when
these intervals intersect. This exact feasibility calculation is distinct
from a likelihood. The equal native up/down record-zero bins cannot
both match the unequal fixed running quark parameters at a common scale
through one such conversion. A physical quark comparison also needs its
renormalization scheme and scale map.

For the separate fixed lepton effects, the sole electron anchor
$0.00051099895069\,\mathrm{GeV}$ gives
\[
 m_\mu^{\mathrm{conditional}}=0.10565941450\,\mathrm{GeV},\qquad
 m_\tau^{\mathrm{conditional}}=1.77698497081\,\mathrm{GeV}.
\]
As a coefficient assignment alone this supplies no propagation proof.
The explicit compiler below supplies a free-field construction on its
declared clock. Anchor-only uncertainty is not a theory error. The
complete target catalogue retains all nine charged fermions, three
neutrino mass eigenstates and five boson types, together with the extra
source modes. Empirical running-action parameters, optional neutrino
coefficients and tree-level bosonic coefficients are distinguished from
source-derived observer rest masses. Global current labels alone supply
no propagating gauge bosons. The exact common-scale test and conditional
GeV table are reproduced by \texttt{model/defect\_mass\_catalogue}
\cite{ophrepo}.

\subsection{Charge-sensitive pole splitting with equal primitive vertices}

Equal primitive quark coefficients need not give equal dressed response
poles. On the unit local source, the total hypercharge $Q_Y$ commutes
with $H$. Its removal charges are $2/3$ for the up-conjugate channel and
$-1/3$ for the down-conjugate channel. A symmetry commuting with $Q_Y$
cannot exchange those two charged operators. The complete spatial
equal-spectrum symmetry instead transforms $Y$ to $-Y+B$, preserving
the baryon-zero singlet domain and reflecting the probe position.
Equality therefore applies to invariant stationary densities and matched
reflected probes; it is not an assertion about an undetermined spatial
ground sector.

A neutral two-state polarization memory with unit occupation cost gives
the declared extension
\[
 H'=H+N_A+Q_YX_A,\qquad
 M_q=\begin{pmatrix}0&q\\q&1\end{pmatrix}.
\]
On a target of charge $q$, every original spectral line splits by
$\lambda_\pm=(1\pm\sqrt{1+4q^2})/2$, with blank-memory weights
$(1\mp1/\sqrt{1+4q^2})/2$. The exact blank resolvent has selfenergy
$\Sigma_q(w)=q^2/(w-1)$. Both branches, every original energy residue,
and preparation failures are retained. The lower dominant source gaps
are approximately $0.4482403502$ and $0.6806484710$ for the up and down
channels. Their ratio $0.6585489711$ fails the fixed common-scale target
comparison. No parameter was chosen using that target ratio.

The unchanged neutral ground has $Q_Y=0$. All nine hundred anomalous
pairing contacts are unchanged, since
$\{[c_i,Q_YX_A],c_j\}=Y_i\{c_i,c_j\}X_A=0$.
Thus unequal propagator poles can arise with equal primitive pairing
vertices. This finite mechanism supplies neither three observed quark
families nor renormalized physical quark masses. Exact controls and
complete records are reproduced by
\texttt{model/quark\_flavour\_resolution} \cite{ophrepo}.

A separate electrostatic extension embeds the full native source as one
charge island on a thirteen-vertex cone, with a compensating reference
and a finite flux register. Solving its Gauss constraints in the positive
electric energy $\sum_e E_e^2/2$ gives $C Q_Y^2$, where
$C=26/217$. Every target energy is translated by $Cq^2$ with its
original residue. These are finite source-reference electrostatic
energies. They supply no transverse photon dynamics, locally gauged
343-site matter source, or intrinsic rest-mass assignment. The exact
Gauss isometry and complete shifted measures are reproduced by
\texttt{model/quark\_gauge\_dressing} \cite{ophrepo}.

\subsection{Neutral composite response and a coupled sterile control}

The native neutral composite
$J_s=\phi_0c_{s,11}-\phi_1c_{s,10}$ commutes with all twelve gauge
currents and carries lepton number $-1$. The original unit source
conserves lepton number. Consequently its anomalous expectation
$\langle J_s(t)J_t(0)\rangle$ vanishes on every invariant reference;
a symmetry-preserving elimination cannot generate a lepton-number-two
term. This does not imply a gapless normal response. The actual local
ground gives nine positive spectral lines. Its normalized cyclic
resolvent has an exact eight-mode complementary selfenergy, with
positive spectral variance approximately $36.4254564105$ and all
original success and failure weights retained.

On the full three-site ninety-fermion-mode source, neutral addition and
removal preparations give an independently executed transport event.
For the matching opposite-spin middle-site receiver,
\[
 \langle R|\psi\rangle=\langle R|H|\psi\rangle=0,\qquad
 \langle R|H^2|\psi\rangle=i/3.
\]
The complete third-order support has $15739$ coordinates. At
$t=10^{-4}$ the spectral remainder bound encloses receiver probability
between approximately $2.5920\times10^{-18}$ and
$2.9700\times10^{-18}$, retaining its full complement and quantum
reference. These preparations are distinct from the local ground.

A declared sterile extension adds the left-handed conjugate $n^c$ of
a right-handed gauge singlet and the unit interaction
\[
 V_\nu=n^c_\uparrow J_\downarrow-n^c_\downarrow J_\uparrow
       +\mathrm{h.c.}
\]
The literal thirty-two-mode algebra verifies all gauge and spin Ward
identities. With quark vertices disabled, the complete eight-fermion-mode
and five-state-scalar control has dimension $1280$, with $80$ physical
Gauss states. Its unique ground energy is $-4$. The active and sterile
neutral probes have positive gaps $1,5$. Sterile removal residues are
$35/72,1/72$ and active-composite residues are $35/72,25/72$;
addition gives the same respective weights. Every other physical energy,
failure and outside-support outcome is retained. The normalized
one-branch sterile propagator is
\[
 G(z)=\frac{z-44/9}{(z-1)(z-5)},\qquad
 \Sigma(z)=\frac{35/81}{z-44/9}.
\]
These two response poles are not counted as two observed neutrino
species. A separate unit sterile Majorana interaction has its own
recomputed ground and complete spectrum. Activating the quark vertices
leaks out of the lepton control, so neither control ground is transferred
to the all-unit full source.

One declared positive energy unit $s$ gives neutral gaps $s,5s$.
For illustration, an empirical normal-ordering scale with zero lightest
mass gives $s=5.01098792655\times10^{-11}\,\mathrm{GeV}$. This unit
choice neither assigns an active or sterile species to either pole nor
establishes a relativistic rest mass. The active tau antineutrino is a
weak-flavour readout, distinct from the sterile coordinate. Three-family
mixing and oscillation observables require their own construction.
The exact source operators, complete spectra and transport bounds are
reproduced by \texttt{model/neutrino\_mass\_response} \cite{ophrepo}.

\subsection{Compiling the lepton effect into a free mass operator}

The literal lepton analyzer acts on $3600$ addresses and three character
coordinates. Executing all $10800$ columns gives a unitary $U$ whose
selected block $K_{00}$ has squared singular values $e_k$ and
$\operatorname{Tr}(K_{00}^\dagger K_{00})=5/18$. Independent extraction
of the configured charged-family polynomial verifies
\[
 F E_{\rm native}F^\dagger=C_{\rm ex}^\dagger C_{\rm ex}/3.
\]
The exact prepared/native-pair factor is $5/4$. Controls reject an extra
square, substitution of the native pair, active-only normalization and
the conjugate Fourier orientation. This identity is a coefficient and
frame theorem, independent of a physical mass assignment.

Lift $U$ to the entire lepton doublet and its conjugate to the singlet
on the full canonical anticommutation algebra. A rank-three selected
projector $Q$ gives $W=U^\dagger Q U$. With neutral background norm one,
vacuum compression of the literal pairing kick gives the coefficient
$Y=K_{00}^\dagger K_{00}$. This is an operator identity on all physical
occupations, without replacing the source Hamiltonian by the desired
effect. There are $90$ physical fermion modes and $324000$ enlarged
modes per site. The mean interaction norm is $5/9$ per site and the
uncompressed kick norm is six per site.

Apply $N$ rounds of duration $t/N$, measure vacuum versus nonvacuum
ancillary occupation, retain the committed classical history and its
quantum outputs, and supply fresh ancillary vacuum. A common free Weyl
kinetic operator commutes with the address analyzer. Interaction-picture
expansion and telescoping give complete unhalved trace-norm error
\[
 \epsilon\le\frac{108 V^2t^2}{N}
\]
for $V$ sites, before separately adding source-dilation errors. Quantum
coherence within each classical history sector is retained. An unmeasured
coherent purification has a different, order $N^{-1/2}$ comparison;
the recorded bound is not applied to it.

The same native Nambu symbol has three charged free Dirac fields per
valley with rest coefficients $e_k$, together with all eight sine
valleys and $39$ zero internal coefficients. A full sixteen-state
canonical transformation identifies the charged pairing with the Dirac
mass Hamiltonian and separately verifies its electric charge. Its
complete addition/removal instrument retains $36$ records and reads
the one-particle energy $e_k$, without an extra square. A compensated
four-effect Ramsey instrument reads probabilities
$(1\pm\cos(e_kt))/4$ and $(1\pm\sin(e_kt))/4$. Complete-output error
$\epsilon$ bounds the recovered energy error by $4\epsilon/t$ on the
specified unaliased phase window. The positive electron denominator is
certified before forming ratios.

One empirical electron anchor gives the three free coefficients
\[
 (0.00051099895069,\ 0.1056594144959,\ 1.7769849708134)\,\mathrm{GeV}.
\]
The fixed phase $2/9$, scalar
background and ancillary preparation are construction inputs. The
uncompressed norm estimate certifies one-site error $\epsilon=10^{-12}$
at $t=1$ with $108\times10^{12}$ rounds; this round count is not executed.
Section~\ref{sec:particle-compiler-compression} gives a sharper bound
and an exact reduction of the same compiler.
Physical clock matching, a selected interacting Higgs vacuum, mirror
removal and observed Standard Model identification are not inferred.
The construction, complete finite controls and error budgets are
reproduced by \texttt{model/lepton\_mass\_compiler} \cite{ophrepo}.

\subsection{Exact active-mode compression and a leakage bound}
\label{sec:particle-compiler-compression}

The same neutral compiler admits an exact reduction of its ancillary
coordinates. For character $k$, let $p$ be the blank address,
$v_k=U_k^\dagger p$, $\lambda_k=\langle p,v_k\rangle>0$ and
$e_k=\lambda_k^2\in(0,1)$. Then
\[
 a_k=\frac{v_k-\lambda_kp}{\sqrt{1-e_k}},\qquad
 B_k=(p,a_k),\qquad
 v_kv_k^\dagger=B_k
 \begin{pmatrix}e_k&\sqrt{e_k(1-e_k)}\\
 \sqrt{e_k(1-e_k)}&1-e_k\end{pmatrix}B_k^\dagger.
\]
The two spin-oriented pairing blocks per character need twelve ancillary
fermion modes in total. Retaining all ninety physical modes gives $102$
modes per site. The exterior lift $\Gamma(T)$ uses $B_k$ on a left field
and $\overline B_k$ on its conjugate right field, fixes the physical
modes, and satisfies
$H_{\rm original}\Gamma(T)=\Gamma(T)H_{\rm compressed}$.
The common Weyl hopping preserves these position-independent spans.
Every omitted address mode starts and stays in vacuum. This identity
holds on the complete physical occupation space and with arbitrary
untouched references. Twelve ancillary modes are minimal among
number-preserving mode embeddings that fix the physical modes and
reproduce this first-order coupling.

The ancillary isometry also intertwines vacuum/nonvacuum flags and
reconstructs the full quantum archive after every round. It preserves
coherence within each committed history. For one oriented pair with kick
$\mathcal V$, let $P$ project onto fresh ancillary vacuum. On its physical
range the leakage form is
\[
 P\mathcal V^2P-(P\mathcal VP)^2
 =(1-e^2)I-e(1-e)(n_L+n_R).
\]
Its maximum is $1-e^2$. Distinct pairs and sites emit into orthogonal
ancillary sectors, so for $V$ sites the full off-diagonal norm obeys
\[
 B^2:=\|(I-P)HP\|^2=V\Bigl(6-2\sum_k e_k^2\Bigr)<6V.
\]
The kinetic term preserves $P$. Block Duhamel integration therefore
gives vacuum-branch operator error at most $B^2h^2/2$ in one round of
duration $h$, independently of the kinetic norm. Telescoping $N$ rounds
and retaining every nonvacuum output gives the complete unhalved bound
\[
 \|\rho_{\rm actual}-\rho_{\rm ideal}\|_1
 \le \min\{2,2B^2t^2/N\}\le\min\{2,12Vt^2/N\}.
\]
Per-round half-diamond implementation error $\eta$ adds $2N\eta$.
The sharper one-site bound at $t=1$ certifies error $10^{-12}$ with
$11724207179903$ rounds; this resource count is not an executed campaign.
The free masses, Ramsey inversion and positive ratio-denominator guard
are unchanged. Complete finite occupation and quantum-archive controls
are reproduced by \texttt{model/lepton\_compiler\_compression}
\cite{ophrepo}.

\subsection{Complete calibrated quadratic spectrum}

A separate reference action uses the explicitly calibrated common-scale
Yukawa, gauge, scalar and optional Weinberg inputs. Its literal native
Higgs tensors give a $45\times45$ fermion mass kernel. Keeping both
momenta $\pm(0,0,1\,\mathrm{GeV})$ retains $180$ physical fermion
modes and the full $360$-dimensional Nambu symbol. An analytic mass frame
preserves every charge, helicity, colour, momentum and conjugate label,
including the small neutral amplitudes. The quadratic dispersion obeys
$E^2-|p|^2=m^2$ in this declared action.

Differentiating the actual gauge-Higgs action gives tree coefficients
approximately
\[
 (80.8480001663,\ 92.2264476377,\ 131.279873449)\,\mathrm{GeV}
\]
for $W$, $Z$ and the scalar, with zero
photon and gluon rest coefficients. A global total-occupation cap two
on the $56$ boson modes has dimension $1653$ and preserves colour
rotations. The complete boson instrument retains $375$ pointer/energy
records, including $362$ zero-weight records. Both fermion and boson
instruments retain their complementary outcomes and compensating
references for the symmetries preserved by the chosen Higgs background.

The catalogue covers seventeen elementary target types, complete
antiparticle and mode multiplicities, and separately calibrated composite
controls. The quark coefficients are supplied running parameters,
the boson coefficients are tree-action quantities, and the optional
neutrino input has its own scale-matching obligations. These conventions
are not merged with the lepton compiler's free mass coefficients or the
quenched hadron fits. Extra sterile, valley, complementary and response
modes are classified as source modes, without declaring undiscovered
particles. The complete finite construction and input ancestry are
reproduced by \texttt{model/conditional\_particle\_spectrum}
\cite{ophrepo}.

\subsection{Joint quantum-number and mass records}

On the same calibrated two-momentum source, electric charge, colour
Casimir, spin about $z$ and momentum have commuting Nambu lifts
\[
 Q_N=\operatorname{diag}(Q,-Q^T),\quad
 C_N=\operatorname{diag}(C,C^T),\quad
 S_N=\operatorname{diag}(S,-S^T),\quad
 P_N=\operatorname{diag}(P,-P^T).
\]
Their commutators with the actual $360$-dimensional Hamiltonian vanish.
At the declared momenta $\pm(0,0,1\,\mathrm{GeV})$, write helicity as
$h=S_NP_N/|P_N|$. Grouping by the measured tuple
$(Q,C,2h,P_z,E)$ gives $84$ fermion projectors of total rank $180$.
Colour-triplet components occupy rank-three projectors. Circular
polarization similarly gives $28$ boson projectors of total rank $56$,
with rank-eight adjoint gluon groups and no longitudinal massless
oscillator. The adjoint Casimir is exactly $3I_8$. Species names enter
only as a comparison dictionary after this grouping.

These $112$ one-excitation projectors, the vacuum and the entire
many-excitation complement sum to identity on the full
$2^{180}\times1653$-dimensional source. A fair mixture of the complete
fermion and boson probes, composed with these $114$ projectors, retains
$570$ pointer/projector records. On the declared vacuum, $113$ are
positive and $457$ are zero. Fermion and boson group probabilities are
respectively rank$/360$ and rank$/116$, and the boson no-operation
vacuum has probability $1/58$. Charge and colour references are
retained. Within a rank-$r$ colour group their conditional output is
coherently entangled; dephasing its colour coordinates changes that
density by unhalved trace norm $2(1-1/r)$.

Exact mass frames attach $E^2=p^2+m^2$ to these projectors. Rational
square-root enclosures distinguish the optional neutrino energies even
when binary64 subtraction would lose their mass differences. At momentum
one GeV, energy bins of half-width $10^{-6}\,\mathrm{GeV}$ merge these
neutral energies, whereas half-width $10^{-26}\,\mathrm{GeV}$ resolves
them in the fixed input packet. This is an arithmetic and instrument
definition, without a physical detector-precision claim. Opposite-helicity
Majorana records do not add sterile species.

Separately, the compiled lepton's literal sixteen-state Dirac block has
fifteen nonzero joint charge/spin/energy projectors. Each family has
sixty complete reservoir and detector records; forgetting spin recovers
the original thirty-six exactly. Momentum is zero and no rest-frame
helicity is assigned. Both constructions are reproduced by
\texttt{model/particle\_joint\_mass\_identity} \cite{ophrepo}.
The calibrated Yukawa packet and electron-anchored compiler retain
their separate input ancestry and clocks.

\subsection{Family invariants and an interference measurement}

The native three-family source has forty-five internal coordinates and
ninety spin-resolved fermion modes. Its twelve gauge currents commute
with all forty-five generators of family-coordinate $U(3)^5$.
With canonical kinetic terms, the charged kernels transform as
\[
 Y_u\mapsto U_Q^\dagger Y_uU_u,\qquad
 Y_d\mapsto U_Q^\dagger Y_dU_d,\qquad
 Y_e\mapsto U_L^\dagger Y_eU_e.
\]
This is equivalence under a change of coordinates when preparations,
detectors and references transform together. The exact generic tangent
rank is $41$ on the $54$ real charged coefficients, leaving thirteen
invariants: nine singular values, three quark mixing angles and one
charge-conjugation/parity (CP) phase. A complex symmetric neutral kernel
transforms as $M_\nu\mapsto U_L^TM_\nu U_L$. Its generic extension has
$66$ real coefficients and orbit rank $44$, leaving twenty-two
invariants. On the rank-two neutral stratum the corresponding count is
$64-44=20$: a zero mass also removes one independent Majorana phase.
The supplied normal-ordering, zero-lightest-mass packet consequently has
nineteen empirical flavour coordinates and one independent unmeasured
Majorana-phase choice.

The native current and Higgs identities admit all twenty-seven complex
charged family coefficients: their embedding has positive diagonal
Gram matrix, and the full ninety-mode pairing is antisymmetric.
They impose no additional linear family restriction. The specified
common-feature trial has $[Y_uY_u^\dagger,Y_dY_d^\dagger]=0$.
This alignment is not forced by those identities. An exact admitted
counterexample uses
\[
 Y_u=\operatorname{diag}(1,2,4),\quad
 Y_d=V\operatorname{diag}(2,3,5),\quad
 Y_e=\operatorname{diag}(1,3,7),
\]
where $V$ has mixing sines $3/5,5/13,8/17$ and phase $\pi/2$. It gives
\[
 J=\operatorname{Im}(V_{00}V_{11}\overline{V_{01}V_{10}})
   =\frac{51840}{830297}.
\]
The actual native weak-raising current has this $V$ block in all three
colour copies of the mass frame. No measured flavour target selects
the counterexample.

A declared two-path current instrument prepares family pairs $(0,1)$
and $(1,0)$ in distinct colour registers and applies $V\otimes V$.
For every output pair $(r,s)$, put $a_{rs}=V_{r0}V_{s1}$ and
$b_{rs}=V_{r1}V_{s0}$. Its eighteen Kraus rows are
\[
 K_{rs,+}=\frac{(a_{rs},-ib_{rs})}{\sqrt2},\qquad
 K_{rs,-}=\frac{(a_{rs},+ib_{rs})}{\sqrt2},\qquad
 \sum_{r,s,\pm}K_{rs,\pm}^\dagger K_{rs,\pm}=I_2.
\]
On the equal path superposition,
$p_{rs,\pm}=|a_{rs}\mp ib_{rs}|^2/4$ and
$J=p_{01,-}-p_{01,+}$. Every other outcome and the path reference are
retained. Conjugating $V$ preserves its singular-value and incoherent
probability data but reverses this contrast. Thus one real signed
contrast supplies information absent from masses and $|V_{ij}|^2$.
The instrument is a coherent current service with a common charge
reference, without a native pulse compilation or laboratory flavour
assignment. Exact orbit calculations, full-source identities and complete
records are reproduced by
\texttt{model/source\_flavour\_identifiability} \cite{ophrepo}.

\subsection{A finite-frequency amputated scalar--fermion response}

The literal $320$-state lepton/scalar control permits a full three-point
calculation with six fermion fields and two scalar fields. Its ground
projector $P_0$ has rank two and energy $-2$. Preparing the marginal
$I_{320}/320$ of a source/dual-reference entangled state and measuring
the complete energy gives $\rho_0=P_0/2$ with probability $1/160$.
All other energy outcomes and the quantum reference are retained.
Define $R(z)=(z-H-2)^{-1}$ and $O=(c_i,c_j,\phi_k^\dagger)$.
For $\omega_1+\omega_2+\omega_3=0$, the complete ordered transform is
\[
 T_{ij,k}=\sum_{(a,b,c)\in S_3}\operatorname{sgn}_F(a,b,c)
 \operatorname{Tr}\!\left[
 \rho_0O_aR(i\omega_a)O_bR(-i\omega_c)O_c\right],
\]
where the sign changes exactly when the fermions exchange order.
This follows by integrating the two positive Euclidean time intervals.
All one-point functions vanish, so the transform is connected. The
full resolvent retains every intermediate state, including $198$
scalar-occupied dark states.

The complete two-point functions used for amputation are
\[
 \begin{split}
 G_{ij}(z)&=\operatorname{Tr}\rho_0
  [c_iR(z)c_j^\dagger-c_j^\dagger R(-z)c_i],\\
 D_{ab}(z)&=\operatorname{Tr}\rho_0
  [\phi_aR(z)\phi_b^\dagger+\phi_b^\dagger R(-z)\phi_a].
 \end{split}
\]
In the native frame, $G$ is diagonal and $D$ is scalar on the doublet.
Equal addition/removal matrices give $D(z)=D(-z)$, fixing the conjugate
scalar-leg convention. At each specified nonzero imaginary frequency,
\[
 \Gamma_{ij,k}(\omega_1,\omega_2)=
 \frac{T_{ij,k}}{G_{ii}(i\omega_1)G_{jj}(i\omega_2)
                         D_{kk}(-i(\omega_1+\omega_2))}.
\]
The exact denominators are nonzero. All seventy-two components at
$(1,2)$ are retained, of which eight are nonzero. In particular,
\[
 \Gamma_{04,0}(1,2)=\frac{474925}{2394532},\qquad
 \Gamma_{04,0}(2,3)=\frac{11640331925}{26140075569}.
\]
These differ from the primitive unit contact. Canonical scalar
normalization $\phi/\sqrt{14/5}$ multiplies $\Gamma$ by $\sqrt{14/5}$.

Summing the four native oriented scalar/pair components and all six
orders at $(1,2)$ gives an operator of norm at most eight. Its quotient
$B$ by eight commutes with all twelve native currents. The four effects
$(I\pm\operatorname{Re}B)/4$ and
$(I\pm\operatorname{Im}B)/4$ are positive and sum to identity.
Their declared square-root instrument gives four successful pointer
records alongside thirteen failed preparation energies, hence seventeen
original records of total weight one. This finite response has no
selected physical mass shell, spacetime momentum or renormalized
Standard Model coupling. Exact full-source resolvents, signs,
propagators and pointer effects are reproduced by
\texttt{model/frequency\_yukawa\_vertex} \cite{ophrepo}.

\subsection{A neutral scalar occupation pole}

On the full thirty-mode local source, let $N$ be the capped scalar
occupation and write the unit-coupling Hamiltonian as $H=V+N$.
The literal operators obey
\[
 N^2=N,\qquad \{2N-I,V\}=0,\qquad H^2-H=V^2,\qquad [N,V^2]=0.
\]
Let $\lambda$ be the $V^2$ eigenvalue selected by the exact ground
projector and put $\Delta=\sqrt{1+4\lambda}$. The two energies in that
invariant pair are $(1\pm\Delta)/2$. The scalar occupation has ground
mean $n=(1-\Delta^{-1})/2$ and connected variance
$w=n(1-n)=\lambda/(1+4\lambda)$. Its retarded response is exactly
\[
 D_N(z)=\frac{2\Delta w}{z^2-\Delta^2},\qquad
 D_q(z)=\frac1{z^2-\Delta^2},\quad
 q=\frac{N-n}{\sqrt{2\Delta w}}.
\]
The positive pole is approximately $23.4611705885$ source units.
The elementary charged scalar field has a different response. The
neutral operator here is constructed from the capped fields themselves:
$\sum_k\{\phi_k^\dagger,\phi_k\}=4I-3N$.

An occupation measurement followed by complete energy measurement has
probabilities $n^2,w$ for occupied records at energy transfers $0,\Delta$
relative to the initial ground, and
$(1-n)^2,w$ for blank records. Each occupation also retains its exact
zero-weight outside-pair outcome. These six records sum to one; the
measurement supplies mean work $2\Delta w$. This is a local scalar
response pole without a spatial propagation or physical Higgs
identification. Exact algebraic intervals and the complete instrument
are reproduced by \texttt{model/native\_scalar\_density\_pole}
\cite{ophrepo}.

\subsection{Pion operators and calibrated composite energies}

The retained quenched Wilson calculation has degenerate light quarks
and no electromagnetic field. Normalized pseudoscalar flavour
matrices $M_+=|u\rangle\langle d|$,
$M_- = |d\rangle\langle u|$ and
$M_0=(|u\rangle\langle u|-|d\rangle\langle d|)/\sqrt2$
have charges $+1,-1,0$ under $[Q,M]$, with
$Q=\operatorname{diag}(2/3,-1/3)$. Identical light-quark propagators give
connected contraction $\operatorname{Tr}(M_a^\dagger M_b)C(t)
=\delta_{ab}C(t)$. The neutral disconnected contraction cancels because
$\operatorname{Tr}M_0=0$. Thus this source supplies one isospin-degenerate
pseudoscalar decay energy for the three pion charge labels.

The archived fit tunes the light-quark point using the supplied ratio
$c=m_\pi/m_\rho=0.1780$. Choosing the empirical rho anchor
$R=0.77526\,\mathrm{GeV}$ then gives
\[
 a^{-1}=\frac{R}{a m_\rho},\qquad
 m_\pi=cR=0.13799628\,\mathrm{GeV},\qquad
 m_N=R\frac{a m_N}{a m_\rho}.
\]
The pion value is fixed by these two inputs. Its conditional fit variance
vanishes by shared-scale cancellation, without implying zero physical
uncertainty. The independently computed nucleon ratio gives approximately
$1.1275\,\mathrm{GeV}$, with paired-stream statistical and fit-sensitivity
errors retained. An alternative conventional string scale
$\sqrt\sigma=0.44\,\mathrm{GeV}$ gives approximately
$0.1054020\,\mathrm{GeV}$ for the degenerate pion. Both conversions are
reported, exposing their scale disagreement.

The complete four-kappa table includes the unresolved lightest nucleon
estimate. This calculation replays reduced correlator and Wilson-loop
tables on $12^3\times24$ at $\beta=5.7$; the raw gauge configurations
are absent. Quark loops, electromagnetic and light-quark splitting,
finite-volume and continuum control, and transfer to the native observer
source and clock are independent physical obligations. The flavour
identities, full fit replay and conditional conversions are reproduced
by \texttt{model/pion\_mass\_readout} \cite{ophrepo}.

\subsection{Correlated hadron fits and paired GeV uncertainties}

The retained thirty-two Wilson streams also support a fit with the full
time covariance in each declared window. If $j_i$ is its leave-one-stream
correlator mean, the compatible stream mean is
$x_i=32\overline j-31j_i$. For $m$ included streams, the covariance of
their mean is
\[
 C=\frac{1}{m(m-1)}\sum_i(x_i-\overline x)(x_i-\overline x)^T.
\]
Every outer jackknife replica refits with its own thirty-one-stream
covariance. The prescribed rules require enough streams, positive
diagonal entries, full correlation rank and condition number at most
$10^8$. All forty-eight windows and all their replicas pass. No
pseudoinverse, shrinkage or discarded replica supplies an accepted fit.
The maximum replica condition number is approximately $4.50\times10^7$.
At fixed decay coefficient $m$, the amplitude is profiled exactly as
$A=f^TC^{-1}y/(f^TC^{-1}f)$; a bounded numerical search over
$0.02\le m\le3$ supplies the fitted coefficient. The fit window and
search rule do not use empirical hadron masses.

Paired stream labels propagate through the chiral extrapolation and
potential-scale conversion. The table gives GeV central values with
statistical errors and the separate variation across the specified
hadron-window, sink, chiral-subset and potential fits:
\[
\begin{array}{cc|rrr}
 \text{calibration}&\text{channel}&m\;(\mathrm{GeV})&
 \sigma_{\rm stat}\;(\mathrm{GeV})&\Delta_{\rm fit}\;(\mathrm{GeV})\\ \hline
 m_\rho=0.77526&\pi&0.13799628&0&0\\
 &\rho&0.77526&0&0\\
 &N&1.130525&0.030231&0.053203\\
 \sqrt\sigma=0.44&\pi&0.106231&0.001401&0.011761\\
 &\rho&0.596801&0.007873&0.066072\\
 &N&0.870287&0.024191&0.076896
\end{array}
\]
The rho-anchored pion and rho values are fixed calibration identities;
their vanishing conditional fit errors do not express physical accuracy.
The correlated nucleon ratio is $1.458252\pm0.038995$. Independent
diagonal-time fits are retained as controls, including the rho-anchored
nucleon value $1.127499\,\mathrm{GeV}$. The inherited lightest nucleon
quality flag is also retained.

This calculation treats the time covariance explicitly; the chiral
regression uses diagonal mass-error weights and is not a fully correlated
errors-in-variables fit. Its errors exclude sea quarks, continuum and
finite-volume extrapolation, anchor uncertainty and native-clock
matching. The same quenched finite source supplies all results in
\texttt{model/correlated\_hadron\_spectrum} \cite{ophrepo}.


\subsection{Open items}

Open are the absolute lepton scale, the quark masses and the split between up and down, the
quark and lepton mixing matrices, the neutrino masses, the Higgs mass term, which is of cell
size in every regulator examined, the declared parts of the coupling chain, and a law that
writes the continuous port phases.
